%% ============================================================
%% NEW SUBSECTION -- "Cost accounting: economic cost versus the solver objective"
%% Insert in the methodology, after the forward evaluator and before the validation
%% protocol. Referenced by the response letter as S2.6.
%%
%% Source: SECTION_2_6_objective_vs_economic.md + objective_decomposition.csv.
%% This section does not exist in the manuscript at all, yet the response letter points
%% R2.2 at it. Without it a reader recomputing the copperplate bias from the objective
%% gets -11.8 % where the paper reports -15.1 %, and has no way to reconcile the two.
%%
%% The three terms the response letter originally named -- return-temperature anchor,
%% terminal storage valuation, demand slack -- are all structurally ZERO on this lineage
%% (verified in the configs and source, not assumed). The residual is CHP CO2 accounting
%% plus storage cycling. That correction is carried here and in the letter.
%% ============================================================

\subsection{Cost accounting: economic cost versus the solver objective}
\label{subsec:costacct}

Every bias and regret figure in this paper is reported on the \emph{economic} operating
cost, not on the raw solver objective, and the two differ systematically. On Memmingen the
objective exceeds the economic cost by $38.7$--$41.0$\,\% depending on the level
(Table~\ref{tab:objdecomp}). Because that difference is large enough to change the
apparent size of every comparison, we state what it is and why the economic cost is the
right quantity.

Two mechanisms account for it, and neither is a distortion of the dispatch. First, carbon
enters the objective on a \emph{gross} basis -- every tonne from fuel combustion and grid
import priced at the carbon price -- whereas the reported cost applies the standard
cogeneration self-use allocation, crediting back the emissions of combined-heat-and-power
electricity that is exported rather than self-consumed. This is the avoided-emissions
convention conventional in district-heating studies, and on this network it amounts to
$55$--$58$\,kEUR. Second, the objective carries a thermal-storage cycling penalty of
$2$\,EUR per MWh of throughput, worth $24$--$26$\,kEUR, which shapes the solution away from
physically implausible rapid cycling but corresponds to no operator cash flow. Together
these reproduce the residual to within about $1.5$\,kEUR per run, under one percent of the
objective.

We report the economic cost because it is what an operator pays and the only quantity a
dispatch decision changes. Both accounting terms are near-constant across fidelity levels
-- they depend on the generation mix and the storage duty, which barely move -- so they
inflate every level's absolute cost by roughly the same amount and therefore \emph{dilute}
the percentage differences between levels rather than distorting them. The copperplate's
estimation bias reads $-11.8$\,\% on the raw objective and $\mathbf{-15.1}$\,\% on the
economic cost: the same finding, measured against a base that does or does not include
terms no operator sees. In every bias and regret \emph{difference} the accounting terms
cancel, so the comparisons themselves are unaffected by the choice.

The carbon convention is itself a modelling decision and we state it explicitly: exported
cogeneration electricity is credited with the emissions it displaces at the grid margin.
The alternative gross convention shifts magnitudes by two to three percentage points --
the copperplate bias reads $-12.2$\,\% rather than $-15.1$\,\% -- and changes neither the
sign of any bias or regret figure nor the loss-dominance result, because the credit is
near-identical at every level. The per-level split is released as
\texttt{objective\_decomposition.csv} so that either basis can be reconstructed.
